% Abstract — English

\thispagestyle{plain}
\begin{center}
  {\Large\bfseries Abstract}
\end{center}
\vspace{0.25cm}

\noindent Population ageing has converted the post-65 phase of life into a long, heterogeneous segment of European societies, yet policy and market responses still treat older adults as a single demographic block defined by age and income. This thesis tests whether the welfare regime that surrounds an over-65 population leaves a measurable imprint on the multidimensional structure of late-life experience itself, contrasting a Mediterranean-familistic context (Italy) with a Nordic-universalistic one (Sweden).

\vspace{0.15cm}

\noindent Using the SHARE Wave~9 release (2021--2022), we build two country-specific analytical samples (2{,}378 Italian and 1{,}787 Swedish respondents aged 65 or older) and segment each independently through a pipeline that combines factor-analytic preprocessing of objective resources (health, wealth, social contact, digital adoption) with subjective dimensions of late-life experience (CASP-12 quality of life, EURO-D depression, loneliness, hope, expected survival) treated as constituent inputs rather than as external validators. K-means partitions the standardised analytical space at $k = 5$ for Italy and $k = 6$ for Sweden; Latent Class Analysis provides cross-method triangulation; the Hungarian assignment algorithm aligns the two partitions in a common feature space for cross-country comparison.

\vspace{0.15cm}

\noindent Four sample-level findings carry the contribution. The \emph{Isolation Paradox}: a 26.9\% Italian segment, labelled \emph{Moderate Isolated}, combines intact health and economic resources with thin social and digital engagement, and shows markedly lower healthcare utilisation than structurally comparable peers; no Swedish profile occupies the analogous position in the standardised space. Fragility takes a binary form in Sweden (a single \emph{Fragile} cluster at 11.5\%) and a graduated form in Italy (\emph{Fragile Resigned} at 13.2\% plus \emph{Fragile Depressed} at 26.7\%, jointly 39.9\% of the analytical sample). Digital adoption saturates the autonomous Swedish profiles (75--99\%) and tracks the resource gradient in Italy without any profile reaching the 75\% threshold. The Swedish high-resource pole rotates with age from a digital-engagement to a wealth-accumulation anchor, against the parallel decline of both Italian high-resource profiles.

\vspace{0.15cm}

\noindent A robustness battery comprising LCA cross-method convergence, specification sensitivity, Wave~8 stability (ARI~$=$~0.16 Italy and 0.23 Sweden), and Hungarian-matched centroid persistence for 4 of 5 Italian and 5 of 6 Swedish profiles supports a structural reading of the four findings as the typological signature of the welfare-regime contrast, with implications for the silver-economy market structure and the design gap on the Italian \emph{Moderate Isolated} segment.

\vspace{0.4cm}

\noindent\textbf{Keywords:} multidimensional segmentation; older adults; SHARE Wave~9; welfare regimes; cross-country comparison; Italy; Sweden; silver economy; subjective wellbeing; K-means; Latent Class Analysis.
